\chapter{Architecture}\label{ch:arch}

This chapter describes how RecaLL is organised at the highest level: the architectural
style, the main building blocks at three levels of detail, the technology chosen for
each layer, the way the AI models are connected, and the build, release and
version-control process. All statements were checked against commit \texttt{60d40f9}
(5~November~2026) of the repository.

\section{Architectural Model/Style Used}\label{sec:arch_style}

RecaLL combines three architectural styles.

\begin{enumerate}[leftmargin=*]
  \item \textbf{Local-first client--server.} The product is a desktop application. A
    Tauri~v2 shell hosts a React~19 single-page interface (the client), and a FastAPI
    process (the server) runs on the same machine at \texttt{127.0.0.1:8000}. In the
    packaged build, the server is a PyInstaller executable that Tauri starts as a
    \emph{sidecar} (\texttt{externalBin: binaries/recall-backend} in
    \texttt{tauri.conf.json}); during development \texttt{run.bat} starts
    \texttt{uvicorn} and the Vite dev server separately. The client and server talk only
    through JSON over HTTP and Server-Sent Events (SSE). All user data (PDFs, extracted
    Markdown, notes, flashcards and review history) stays in local SQLite files and a
    local folder. The only outbound traffic is to the AI provider the student chooses,
    and to Supabase if the optional cloud sync is configured.
  \item \textbf{Layered server.} Inside the backend, code is arranged in strict layers:
    HTTP routers $\rightarrow$ services $\rightarrow$ repositories $\rightarrow$ SQLite.
    Routers (\texttt{app/api/v1/*.py} and the study routes in \texttt{app/main.py})
    handle validation, authentication and response envelopes; services
    (\texttt{app/services/}, \texttt{app/socratic\_drill.py},
    \texttt{app/llm\_segment.py}) hold the business rules; repositories
    (\texttt{app/models/repositories.py} and \texttt{app/database.py}) hold all SQL.
  \item \textbf{Plug-in (strategy) layer for variable parts.} The two parts that vary
    most, the AI provider and the PDF extractor, are hidden behind abstract interfaces
    (\texttt{BaseLLMProvider}, \texttt{BaseAIAdapter} and \texttt{BaseExtractor}) and
    chosen at run time by a factory or a registry. New providers or extractors can be
    added without changing the features that use them.
\end{enumerate}

The document-processing path is additionally organised as a \emph{pipe-and-filter}
pipeline (validate $\rightarrow$ extract $\rightarrow$ normalise $\rightarrow$ chunk
$\rightarrow$ embed $\rightarrow$ index), which is described in \Cref{sec:arch_hld}.

\subsection{Rationale for Choosing the Architectural Model/Style}\label{sec:arch_rationale}

The style was chosen to serve the quality attributes that matter most for a study tool
that holds a student's own course material~\cite{bass2021}:

\begin{itemize}[leftmargin=*]
  \item \textbf{Privacy and ownership.} Lecture slides and personal notes never have to
    be uploaded to a server run by the team. Only the text of the topic being processed
    is sent to the AI provider that the student selected, using the student's own key
    (BYOK), or nothing at all when a local Ollama model is used.
  \item \textbf{Offline use and low cost.} With Ollama the whole system runs without an
    Internet connection. There is no hosted backend to pay for, which matters for a
    student project with no budget.
  \item \textbf{Maintainability.} The layered backend keeps SQL out of the routers and
    HTTP out of the services, so each part can be tested on its own. The test suite
    swaps the database path (\texttt{set\_db\_path}) and the AI adapter
    (\texttt{MockAIAdapter}) without touching feature code.
  \item \textbf{Modifiability of the AI part.} LLM vendors change models, prices and
    quotas often (the code already contains a fallback from a quota-limited Gemini
    model to \texttt{gemini-flash-latest}). The strategy layer confines such changes to
    one provider class.
  \item \textbf{Performance.} The heavy work (PDF parsing, chunking) runs next to the
    files, without network upload. FastAPI's \texttt{asyncio} model lets one process
    stream LLM output to the interface while other requests are served.
\end{itemize}

\paragraph{Alternatives considered and rejected.}
\begin{itemize}[leftmargin=*]
  \item \emph{Cloud SaaS (central server and database).} It would make multi-device use
    easy, but it would require hosting costs, server-side storage of copyrighted course
    PDFs, and a stronger security and privacy burden. Parts of a SaaS design were kept
    (accounts, workspaces, plans in \texttt{recall\_saas.db}) so that a hosted version
    remains possible, but they run locally.
  \item \emph{Microservices.} Separate services for ingestion, AI and review would add
    network hops, deployment work and distributed failure modes that a four-person
    team could not run in one term, and a desktop install would need several
    processes.
  \item \emph{Electron instead of Tauri.} Electron bundles a full Chromium. Tauri uses the
    operating system's WebView and a Rust core, which gives a much smaller installer and
    a first-class sidecar mechanism for the Python backend.
  \item \emph{All logic in the client (no Python backend).} The best PDF tools
    (PyMuPDF, PyMuPDF4LLM, Marker) and the FSRS reference library are Python packages,
    so a Python process was the practical choice.
\end{itemize}

\section{Architectural Design (High-Level Design)}\label{sec:arch_hld}

The architecture is shown at three levels of detail. Level~0 shows the system and its
external actors, level~1 the layers and main services, and level~2 the two most
complex subsystems: the document-processing pipeline and the AI provider layer.

\paragraph{Level 0: system context.}
\Cref{fig:arch_l0} shows the whole system as six blocks: the Tauri desktop app with
its React user interface, the FastAPI route layer, the service layer, local storage
(two SQLite databases and a file folder), the AI provider layer, and the external LLM
services. Only the last block is outside the student's machine, and it is optional: with
a local Ollama server the system runs fully offline.

\figplaceholder[0.85\textwidth]{arch_l0.pdf}{Architecture level 0: RecaLL in its environment}{fig:arch_l0}

\paragraph{Level 1: layers and services.}
\Cref{fig:arch} opens the desktop application into its layers. The UI layer is the
React interface inside the Tauri shell. The API layer has two route groups: the
versioned \texttt{/api/v1} routers, which use JWT bearer authentication and the
\texttt{recall\_saas.db} database, and the older study routes in
\texttt{app/main.py} (books, topics, notes, annotations, Socratic drill, settings and
prompts), which use a per-user \texttt{user\_<id>\_recall.db}. Both groups call the
same service, provider and storage layers.

\figplaceholder[0.95\textwidth]{arch_l1.pdf}{Architecture level 1: layers, route groups and services}{fig:arch}

\paragraph{Level 2: document-processing pipeline.}
\Cref{fig:arch_l2_pipeline} shows the stages that turn an uploaded PDF into searchable
chunks. In the \texttt{/api/v1/documents/upload} path, the upload handler stores the
file, creates a \texttt{processing\_jobs} row and schedules
\texttt{DocumentProcessingPipeline.process\_document} as a FastAPI background task. Each
stage writes its name and a progress value (10, 30, 45, 60, 80, 95 and finally 100
per cent) to the job row, which the interface polls through
\texttt{/documents/\{id\}/status}. The study path (\texttt{/books/upload} followed by
\texttt{/books/\{id\}/process-stream}) uses the extractor registry, builds a topic tree
from the PDF's table of contents, and streams progress as SSE.

In the shipped interface (\texttt{UploadDocumentModal.tsx}) a PDF is uploaded through
the study path, \texttt{POST /books/upload}, so that the reader, topic tree, notes and
drill work on it. The modal then also sends the same file to
\texttt{/api/v1/documents/upload} as a background request. In
addition, \texttt{BookSyncService} (\texttt{app/services/book\_sync.py}) mirrors
imported books into \texttt{recall\_saas.db}, chunking them with the same
\texttt{SemanticChunker} (1\,500/150), so the Knowledge Hub can search them. Markdown
and text files go only through \texttt{/api/v1/documents/upload}.

The two paths use different extractors. The \texttt{/api/v1} pipeline calls
\texttt{ExtractionService} (\texttt{app/services/extractor.py}), which checks size,
extension and the \texttt{\%PDF} signature and reads PDF, Markdown and plain-text files
with PyMuPDF (\texttt{fitz}) directly. The extractor plug-in registry (PyMuPDF4LLM by
default, Marker optional) is used only by the study path (\texttt{pdf\_extract.py},
\texttt{toc\_parser.py}, \texttt{topic\_subdivider.py}); on this path scanned or
handwritten PDFs are read with Marker, installed locally or through the Datalab API. In
the embedding stage, \texttt{EmbeddingService} calls an OpenAI, Gemini or Ollama
embedding API when a server key or Ollama host is configured in \texttt{.env}.

\figplaceholder[0.95\textwidth]{arch_l2_pipeline.pdf}{Architecture level 2: document-processing pipeline}{fig:arch_l2_pipeline}

\paragraph{Level 2: AI provider layer.}
\Cref{fig:arch_l2_ai} shows that there are two AI stacks, one per route group.
\begin{itemize}[leftmargin=*]
  \item \emph{Study stack} (\texttt{app/llm\_providers}). The legacy study features
    (segmentation, Cornell notes, Socratic drill, topic chat, explanations) fill a
    template from the prompt registry, ask \texttt{get\_llm\_provider()} for a provider,
    and call \texttt{generate()} or \texttt{generate\_stream()}. The providers retry
    transient errors with exponential back-off; the Gemini provider also passes a
    priority token-bucket rate limiter and reports to Langfuse. The raw text is
    repaired, parsed and validated against a Pydantic model.
  \item \emph{v1 stack} (\texttt{app/services/ai}). The \texttt{/api/v1} services
    (Knowledge Hub, quizzes, flashcard sets, summaries and concepts) build their prompts
    inline and call \texttt{AIService}, which picks the provider from the user's
    preference, then a valid BYOK credential, then the server setting; decrypts the
    BYOK key (AES-256-GCM); builds an OpenAI-compatible, Gemini or Ollama adapter; and
    falls back to a deterministic \texttt{MockAIAdapter} when no key exists. These
    adapters have no retry or rate limiter.
\end{itemize}

\figplaceholder[0.95\textwidth]{arch_l2_ai.pdf}{Architecture level 2: AI provider layer}{fig:arch_l2_ai}

\paragraph{Component descriptions.}
\Cref{tab:components} summarises every component in the level-1 and level-2 views.

{\small
\begin{longtable}{@{}p{0.27\textwidth}p{0.68\textwidth}@{}}
\caption{Main components of RecaLL and their responsibilities}\label{tab:components}\\
\toprule
\textbf{Component} & \textbf{Responsibility}\\
\midrule
\endfirsthead
\toprule
\textbf{Component} & \textbf{Responsibility}\\
\midrule
\endhead
\bottomrule
\endfoot
Tauri shell (\texttt{desktop/src-tauri}) & Native window, deep links (\texttt{recallai://}), secure key storage plug-ins, and start/stop of the backend sidecar.\\
React UI (\texttt{desktop/src}) & About 20 views (Dashboard, Documents, Knowledge Hub, Notes, Review, Flashcards, Quizzes, Progress, Settings, LLM Inspection, \ldots), Redux Toolkit slices for shared state, and a typed API client (\texttt{src/api/client.ts}).\\
\texttt{/api/v1} routers & 14 router modules, 87 endpoints: auth, users, health, documents, search, conversations, flashcards, quizzes, attempts, learning, reviews, knowledge, billing and account.\\
Study routes (\texttt{app/main.py}) & 71 endpoints for books, topics, Cornell notes, annotations, the Socratic drill, legacy flashcards, chat, settings, prompts and token usage.\\
Middleware (\texttt{app/api/middleware.py}) & Security headers, request-size limits (10\,MB, 55\,MB for uploads), CORS, and an in-memory login rate limiter.\\
Dependencies (\texttt{app/api/deps.py}) & \texttt{get\_current\_user}, \texttt{get\_current\_workspace}, \texttt{require\_document\_owner} and \texttt{get\_current\_user\_optional}.\\
Document pipeline (\texttt{app/services/\allowbreak{}pipeline.py}) & Seven-stage processing with job tracking and failure codes.\\
Extraction & v1: \texttt{ExtractionService} (PyMuPDF; PDF, Markdown, text). Study path: registry in \texttt{app/extractors} with plug-ins \texttt{pymupdf4llm} (default), \texttt{marker} and \texttt{marker\_api}; only those whose libraries are installed are offered.\\
Chunker and embeddings & \texttt{SemanticChunker} (1\,500 characters, 150 overlap) and \texttt{EmbeddingService} (OpenAI, Gemini or Ollama embeddings).\\
Retrieval and RAG (\texttt{retrieval.py}, \texttt{rag.py}) & Hybrid keyword/vector scoring, neighbour suppression, a 3\,000-token context builder, grounded answers with citation checking.\\
Learning services & Flashcards, quizzes and attempts, review queue and sessions, the FSRS scheduler, learning analytics, summaries and concepts.\\
Study engine & \texttt{llm\_segment.py} (topic segmentation, flashcards, summaries, explanations, chat), \texttt{socratic\_drill.py} (questions and grading), and Cornell note generation.\\
Prompt registry (\texttt{app/prompt\_manager.py}) & Eleven named templates with declared variables; user overrides stored in the \texttt{system\_prompts} table.\\
LLM providers (\texttt{app/llm\_providers}) & \texttt{BaseLLMProvider}, a factory, Gemini, OpenAI (also used for OpenRouter), Groq and Ollama providers, a rate limiter and a retry helper.\\
AI service (\texttt{app/services/ai}) & Adapter layer used by \texttt{/api/v1}: resolves the provider from preferences, BYOK vault or server settings, and falls back to a mock adapter.\\
Security (\texttt{app/core/security.py}) & PBKDF2-HMAC-SHA256 (600\,000 rounds) passwords, HS256 JWTs, and AES-256-GCM encryption of v1 BYOK keys.\\
Storage & \texttt{recall\_saas.db} (29 tables including \texttt{schema\_migrations}), \texttt{user\_<id>\_recall.db} (14 tables), and the local file store for PDFs and extracted images.\\
Optional sync (\texttt{app/sync\_service.py}) & Push/pull of study data and PDFs to Supabase; skipped when Supabase is not configured.\\
\end{longtable}
}

\subsection{Data-Flow Diagrams (Levels 0, 1 and 2)}\label{sec:dfd}

The architecture views above show \emph{where} code runs; the data-flow diagrams show
\emph{what data} moves between the student, the processes and the stores. They are
levelled in the usual way: the context diagram (level~0) shows the whole system as one
process, level~1 splits it into five numbered processes, and level~2 decomposes the
three processes that matter most for learning (2.0, 3.0 and 4.0) into numbered
sub-processes (2.1, 2.2, \ldots). Flows that enter or leave a parent process appear
again on its child diagram, so the levels balance. The diagrams use Gane--Sarson
notation: rounded rectangles are processes, open-ended rectangles with an ID cell are
data stores, and squares are external entities. A data store that is drawn twice on one
diagram gets an extra bar in its ID cell, and an external entity drawn twice gets a
tick in its corner.

\widefig{dfd_level0.pdf}{Data-flow diagram, level 0 (context diagram)}{fig:dfd0}

\paragraph{Level 0.} \Cref{fig:dfd0} shows RecaLL as the single process 0. The student
sends PDFs, questions, typed answers and review ratings, and receives answers,
flashcards, quizzes, scores and progress. The AI provider (Gemini, OpenAI, Groq,
OpenRouter or a local Ollama model) receives prompts with document text and returns
generated text or JSON. Three external entities are optional: Supabase (login token and
cloud sync), the mock billing provider (checkout request and webhook), and the Marker or
Datalab OCR service, which turns scanned pages into Markdown.

\widefig{dfd_level1.pdf}{Data-flow diagram, level 1}{fig:dfd1}

\paragraph{Level 1.} \Cref{fig:dfd1} opens process 0 into five processes: 1.0
Authenticate and manage plan, 2.0 Ingest document, 3.0 Answer questions, 4.0 Review
cards (FSRS) and 5.0 Generate study material. They share five data stores: D1 Accounts
(users, plans, usage), D2 Documents (files, chunks), D3 Conversations (messages,
sources), D4 Legacy study (books, topics, notes, cards and drills in the per-user
database) and D5 Study records (flashcard sets, quizzes, review sessions). D1, D2, D3 and
D5 live in \texttt{recall\_saas.db}; D4 is the per-user study database.

\figplaceholder[0.95\textwidth]{dfd_level2_process_doc.pdf}{Data-flow diagram, level 2 of process 2.0 Ingest document}{fig:dfd2-doc}

\paragraph{Level 2: 2.0 Ingest document.} \Cref{fig:dfd2-doc} follows a file from the
student to indexed chunks. The upload first checks the size limits and the plan quota
against D1 (2.1), then saves the file and a processing job in D2 (2.2) and registers the
book and its topics in the legacy store D4 (2.3). The background job then validates the
file and extracts page text (2.4), splits it into chunks (2.5), obtains embeddings from
the AI provider or the placeholder generator (2.6) and indexes the chunks with their
vectors in D2 (2.7). The student's interface polls the job status (2.8) until the
document is ready or has failed.

\figplaceholder[0.95\textwidth]{dfd_level2_ask.pdf}{Data-flow diagram, level 2 of process 3.0 Answer questions}{fig:dfd2-ask}

\paragraph{Level 2: 3.0 Answer questions.} \Cref{fig:dfd2-ask} is the Knowledge Hub.
A question is accepted only after access and credits are checked in D1 (3.1); it is
saved to D3 (3.2); the most relevant chunks are retrieved from D2 (3.3); and a grounded
prompt is built from the chunks and the recent history (3.4). The AI provider generates
the answer (3.5), citation markers that do not match a retrieved chunk are removed
(3.6), and the cited answer and the token and credit usage are stored in D3 and D1
(3.7, 3.8).

\figplaceholder[0.95\textwidth]{dfd_level2_review.pdf}{Data-flow diagram, level 2 of process 4.0 Review cards}{fig:dfd2-review}

\paragraph{Level 2: 4.0 Review cards.} \Cref{fig:dfd2-review} shows the spaced-repetition
loop on the D5 stores. Due items are found in D5.1 Learning items (4.1) and copied into
a review session in D5.2 (4.2, at most 20 items). Each card is shown front first and
then revealed from D5.3 (4.3). The student's rating is passed to the FSRS scheduler
(4.4), which returns the new memory state and due time; the learning item is updated in
D5.1 (4.5) and an immutable review event is written to D5.4 (4.6).

\section{Technology, Software, and Hardware Used}\label{sec:tech}

\Cref{tab:stack} lists the technology of each layer. Versions are the ones installed in
the development environment at the final commit.

{\small
\begin{longtable}{@{}p{0.16\textwidth}p{0.27\textwidth}p{0.50\textwidth}@{}}
\caption{Technology stack and justification}\label{tab:stack}\\
\toprule
\textbf{Layer} & \textbf{Technology} & \textbf{Justification}\\
\midrule
\endfirsthead
\toprule
\textbf{Layer} & \textbf{Technology} & \textbf{Justification}\\
\midrule
\endhead
\bottomrule
\endfoot
Desktop shell & Tauri 2.11 (Rust) & Small installer using the system WebView; built-in sidecar support for the Python backend; plug-ins for deep links, file system and secure storage~\cite{tauri_docs}.\\
Frontend & React 19.2, TypeScript 6 (strict), Vite 8 & Component model fits the many study views; strict typing catches API-contract errors at build time; Vite gives fast reloads and builds.\\
UI and state & Tailwind CSS 4, React Router 7, Redux Toolkit 2, Framer Motion & Utility CSS keeps the neo-brutalist design consistent; Redux holds reader, review and provider state shared across views.\\
Rich content & KaTeX 0.18, react-markdown, pdf.js (pdfjs-dist 6), BlockNote 0.52, Recharts 3 & Mathematics in notes and flashcards; in-app PDF reading and highlighting; a Notion-style note editor; progress charts.\\
Backend & Python 3.12 (CI) / 3.13 (local), FastAPI 0.162, Uvicorn 0.54, Pydantic 2.14 & Async I/O for streaming LLM output; automatic OpenAPI schema; request and LLM-output validation with Pydantic~\cite{fastapi_docs}.\\
PDF processing & PyMuPDF 2.1, PyMuPDF4LLM 2.1, Marker (optional) & Fast text, TOC and image extraction; PyMuPDF4LLM converts pages to Markdown for the study path, and Marker adds OCR for scanned and handwritten pages there~\cite{pymupdf4llm}.\\
Scheduling & \texttt{fsrs} 6.3.2 (FSRS-6) & Reference implementation of the Free Spaced Repetition Scheduler, which models memory stability and difficulty per card~\cite{ye2022fsrs,pyfsrs}.\\
AI & Gemini, OpenAI, Groq, OpenRouter (HTTP via httpx 0.28), Ollama (local) & No model is trained; hosted or local LLMs are called through one interface so students can pick by cost, quota or privacy. \texttt{json-repair} 0.63.3 fixes malformed JSON.\\
Database & SQLite (WAL mode) & Zero-administration embedded database that ships with Python; WAL allows reads during writes~\cite{sqlite_wal}. \texttt{docker-compose.yml} and \texttt{migrations/001\_initial\_schema.sql} describe a PostgreSQL + pgvector variant prepared for a future hosted version; the desktop product uses the \texttt{sqlite3} driver.\\
Security & PyJWT 2.15, \texttt{cryptography} 50 (AES-GCM), \texttt{hashlib} PBKDF2 & Standard, audited primitives; no custom cryptography.\\
Observability & Langfuse 5.1 (optional) & Traces of LLM calls and errors when keys are configured; a no-op otherwise.\\
Cloud sync & Supabase 2.33 (optional) & Authentication and storage for multi-device sync, used when Supabase keys are configured.\\
Testing and quality & pytest 9.2, Playwright 1.62, oxlint 1.74, \texttt{tsc} & Backend unit and integration tests, browser end-to-end journeys, fast linting and type checking.\\
CI/CD & GitHub Actions, PyInstaller, \texttt{tauri-action} & Automated test, build and release on GitHub~\cite{gha_docs}.\\
Hardware & Four Windows developer laptops, each with a 512~GB SSD: Intel Core~i5 with 32~GB RAM, AMD Ryzen~5 with 16~GB, AMD Ryzen~7 with 16~GB and Intel Core~i5 with 16~GB (\Cref{tab:hw}) & The app targets an ordinary student laptop; heavy work is done by the remote LLM, or by a small local model (default \texttt{gemma3:4b} in the core settings) when Ollama is used.\\
\end{longtable}
}

\section{AI-integration Architecture}\label{sec:ai_arch}

RecaLL does not train or host its own model. It integrates pre-trained large language
models as a service, and wraps every call in a fixed sequence of steps
(\Cref{fig:aiarch}).

\figplaceholder[0.95\textwidth]{ai_architecture.pdf}{AI-integration architecture: prompt, context, provider, validation and fallback}{fig:aiarch}

\begin{enumerate}[leftmargin=*]
  \item \textbf{Data pipeline.} Source text comes from the document pipeline: the
    Markdown of a topic (study path) or the stored chunks of a document (\texttt{/api/v1}
    path). The Socratic drill strips image links from the text first, while Cornell
    notes pass the list of extracted figures so they can be embedded. Every input is
    cut to a budget (8\,000 characters for Cornell notes and drills, 7\,000
    characters for a single concept, 3\,000 estimated tokens of retrieved context for the
    Knowledge Hub).
  \item \textbf{Prompts.} Each study task has a named template in
    \texttt{PROMPT\_REGISTRY} with a list of required variables. A student may override
    any template through \texttt{/api/prompts}; the override is checked for missing or
    undefined placeholders before it is saved. The \texttt{/api/v1} services do not use
    the registry; their prompts (for example the Knowledge Hub system prompt in
    \texttt{rag.py}) are built inline.
  \item \textbf{Provider resolution} (two stacks, \Cref{fig:arch_l2_ai}). The study features call
    \texttt{get\_llm\_provider()}, which reads the saved provider type and model from
    the settings table. The \texttt{/api/v1} features call
    \texttt{AIService.resolve\_adapter()}, which tries, in order, the user's preferred
    provider, a valid BYOK credential, the server settings, and finally a deterministic
    \texttt{MockAIAdapter}.
  \item \textbf{Inference.} All models are reached over REST (HTTPS for hosted models,
    \texttt{http://localhost:11434} for Ollama). JSON-producing tasks pass a JSON schema
    derived from the Pydantic model; OpenAI-compatible providers use
    \texttt{response\_format} with a strict schema, Gemini uses
    \texttt{responseMimeType: application/json}, and Ollama uses its \texttt{format}
    field.
  \item \textbf{Validation and repair.} Responses are cleaned (code fences, bad escapes,
    trailing commas, truncated brackets), parsed, and validated with Pydantic. The
    Socratic drill retries once with a stricter instruction if parsing fails.
  \item \textbf{Resilience and fallback.} In the study stack, transient HTTP errors
    (429, 502, 503, 504) and time-outs are retried with exponential back-off and jitter;
    Gemini calls also pass a priority-aware token-bucket rate limiter; and a
    quota-exhausted Gemini model falls back to \texttt{gemini-flash-latest}. The v1
    adapters do not retry; they convert provider errors into typed exceptions
    (authentication, rate limit with \texttt{Retry-After}, time-out, unavailable) that
    the routes return to the user. When no key is configured at all, the
    \texttt{/api/v1} features use the deterministic mock adapter.
  \item \textbf{Monitoring.} The study-stack providers write token usage per provider,
    model and feature to \texttt{token\_usage\_logs}, which the LLM Inspection page
    shows; the v1 services record credit usage in \texttt{usage\_records}. When Langfuse
    keys are set, study-stack generations and errors are traced there as well.
\end{enumerate}

The design of the prompts and of the validation step is detailed in
\Cref{sec:ai_design}, and the code in \Cref{sec:ai_impl}.

\section{CI/CD Pipeline}\label{sec:cicd}

The repository has two GitHub Actions workflows (\Cref{fig:cicd}). They are real and run
on GitHub; they are not a plan.

\figplaceholder[0.95\textwidth]{cicd_pipeline.pdf}{CI and release pipeline stages}{fig:cicd}

\paragraph{Continuous integration (\texttt{ci.yml}).}
The workflow runs on every push and pull request to \texttt{main} (or
\texttt{master}) and can be started by hand. It has two independent jobs on
\texttt{ubuntu-latest}:
\begin{itemize}[leftmargin=*]
  \item \textbf{Backend Pytest Suite:} checkout $\rightarrow$ Python 3.12 with pip cache
    $\rightarrow$ install \texttt{requirements.txt} plus pytest tools $\rightarrow$
    \texttt{python -m pytest app/tests -v --durations=10} with \texttt{TESTING=true}.
  \item \textbf{Frontend Lint, Build \& E2E:} checkout $\rightarrow$ pnpm~9 and
    Node~20 with a pnpm cache $\rightarrow$ \texttt{pnpm install --frozen-lockfile}
    $\rightarrow$ \texttt{pnpm run lint} (oxlint) $\rightarrow$ \texttt{pnpm run build}
    (\texttt{tsc -b \&\& vite build}) $\rightarrow$ install Chromium for Playwright
    $\rightarrow$ start \texttt{vite preview} on port 4173 $\rightarrow$
    \texttt{pnpm run test:e2e}.
\end{itemize}
\Cref{lst:ci} shows the backend job and the first steps of the frontend job, copied
from the repository.

\begin{lstlisting}[caption={Excerpt of \texttt{.github/workflows/ci.yml}},label={lst:ci}]
name: CI
on:
  push:
    branches: [main, master]
  pull_request:
    branches: [main, master]
  workflow_dispatch:

jobs:
  backend-tests:
    name: Backend Pytest Suite
    runs-on: ubuntu-latest
    steps:
      - name: Checkout repository
        uses: actions/checkout@v4

      - name: Set up Python 3.12
        uses: actions/setup-python@v5
        with:
          python-version: '3.12'
          cache: 'pip'

      - name: Install Python dependencies
        run: |
          python -m pip install --upgrade pip
          pip install -r requirements.txt
          pip install pytest pytest-cov pytest-asyncio httpx

      - name: Run Pytest Test Suite
        env:
          TESTING: 'true'
          DATABASE_URL: 'sqlite:///:memory:'
        run: |
          python -m pytest app/tests -v --durations=10

  frontend-tests:
    name: Frontend Lint, Build & E2E
    runs-on: ubuntu-latest
    steps:
      ...
      - name: Run Oxlint Linter
        working-directory: ./desktop
        run: pnpm run lint

      - name: Run TypeScript & Vite Build
        working-directory: ./desktop
        run: pnpm run build
\end{lstlisting}

\paragraph{Release (\texttt{release.yml}).}
Pushing a tag that matches \texttt{v*} (or a manual dispatch) starts the release job on
\texttt{windows-latest} with \texttt{contents: write} permission. It sets up
Python~3.12, installs the requirements and PyInstaller, and runs
\texttt{scripts/build\_sidecar.py}, which freezes the FastAPI backend into
\texttt{desktop/src-tauri/binaries/recall-backend-<target-triple>.exe}. It then sets up
pnpm~9 and Node~20, builds the frontend, installs the stable Rust toolchain for
\texttt{x86\_64-pc-windows-msvc}, and runs \texttt{tauri-apps/tauri-action@v0}, which
compiles the Tauri app, bundles the sidecar and produces the Windows installers. The
workflow passes no \texttt{tagName} or \texttt{releaseName} input to the action, so
\texttt{tauri-action} only builds the installers; it does not create or publish a GitHub
Release.

\section{Git Management (Version Controlling)}\label{sec:git}

\begin{itemize}[leftmargin=*]
  \item \textbf{Repository:} \url{https://github.com/Saimun-jd/RecallAI}. The submitted
    commit is \texttt{60d40f9} on \texttt{main} (5~November~2026). The first commit,
    \texttt{cf9bb3f} ``initial commit'', is dated 22~July~2026.
  \item \textbf{Branching strategy:} a lightweight trunk-based flow
    (\Cref{fig:git_workflow}). \texttt{main} is the only long-lived branch and most
    commits were pushed to it directly. Short-lived branches were used for automated
    fixes: the remote holds three \texttt{copilot/*} branches created by the GitHub
    Copilot agent, including \texttt{copilot/fix-release-windows-latest}. \texttt{main}
    has no branch protection and no required reviews on GitHub.
  \item \textbf{Pull requests and review:} pull request \#1 (``Fix pnpm workspace config
    for release build'', from \texttt{copilot/fix-release-windows-latest}) was merged by
    Al Rubayet on 5~November~2026 as merge commit \texttt{0df9f85}; a second Copilot pull
    request, \#2, is still open and has not been merged. All other changes
    were pushed to \texttt{main} without a pull request; one local merge
    (\texttt{ad823bf}) records parallel work on \texttt{main}. Review inside the team
    was informal: members checked each other's work at the weekly in-person meetings, and
    apart from PR~\#1 there were no formal pull-request reviews.
  \item \textbf{CI history:} \texttt{release.yml} existed by the time of the
    \texttt{v0.1.5} tag (August); \texttt{ci.yml} was added on 25~September~2026 in
    commit \texttt{5ec7f5a}.
  \item \textbf{Tags and releases:} a tag matching \texttt{v*} triggers
    \texttt{release.yml}. HEAD is not tagged. The only tag is \texttt{v0.1.5} on commit \texttt{24667c8}
    (24~August~2026), which matches the Tauri \texttt{productName} ``Recall'' version
    0.1.5.
  \item \textbf{Commit messages:} the later history uses Conventional Commits
    (\texttt{type(scope): summary}), for example \texttt{fix(sync): gracefully skip
    download\_missing\_pdfs\_stream when Supabase is unconfigured} and
    \texttt{feat: add review queue service and interactive review UI components}.
    17 of the 48 commits follow the convention strictly (8 \texttt{feat}, 6 \texttt{fix},
    2 \texttt{chore}, 1 \texttt{docs}); the earlier commits use free-form messages such as
    ``first exe published''.
  \item \textbf{Ignored files and secrets:} \texttt{.env} is ignored and
    \texttt{.env.example} documents the variables; commit \texttt{9fd02fc}
    (``chore: remove hardcoded api keys'') removed keys that had been committed early
    on; those keys were revoked.
\end{itemize}

\figplaceholder[0.9\textwidth]{git_workflow.pdf}{Git workflow: \texttt{main}, short-lived branches and pull requests, tags that trigger releases}{fig:git_workflow}

\paragraph{Contribution statistics.}
\Cref{tab:git_stats} summarises the history from \texttt{git log} and
\texttt{git shortlog -sne}. Commit counts measure activity, not effort; for example a
single commit (\texttt{5ec7f5a}) delivered a whole hardening increment. The individual
contributions are discussed in \Cref{ch:pm}.

\begin{table}[htbp]
\centering\small
\caption{Repository statistics at commit \texttt{60d40f9}}\label{tab:git_stats}
\begin{tabularx}{\textwidth}{@{}lX@{}}
\toprule
\textbf{Measure} & \textbf{Value}\\
\midrule
Commits on \texttt{main} & 48 (2 of them merges)\\
Commits by author identity & Saimun-jd 31; rubayet2027 14; copilot-swe-agent[bot] 2; ``Al Rubayet'' (GitHub web merge by rubayet2027) 1\\
Mapping of identities to team members & \texttt{Saimun-jd} is Saimun (Imteajul Islam); \texttt{rubayet2027} and ``Al Rubayet'' are Turjo (Al Rubayet). Most \texttt{rubayet2027} commits were made on Saimun's computer while it was logged in to Turjo's account. Tahsin and Rony have no commits under their own names\\
Commits per month & July 9; August 16; September 19; November 4\\
Lines changed (all commits) & 113\,431 insertions, 16\,712 deletions\\
Tracked files & 400 (151 Python files; 181 TypeScript/TSX files in \texttt{desktop/src})\\
Size of main source trees & \texttt{app/}: about 38\,300 lines of Python; \texttt{desktop/src}: about 37\,100 lines of TypeScript/TSX\\
Branches on the remote & \texttt{main} and three \texttt{copilot/*} branches\\
Tags & \texttt{v0.1.5}\\
\bottomrule
\end{tabularx}
\end{table}
